\documentclass{article}
\usepackage[utf8]{inputenc}
\usepackage[french]{babel}
\usepackage{graphicx}
\usepackage{hyperref}

\title{Rapport de Projet : Medical Agentic Research Assistant}
\author{Manus AI}
\date{\today}

\begin{document}

\maketitle

\section{Introduction}
Ce document présente le développement d'un assistant de recherche médicale intelligent basé sur une architecture agentique.

\section{Architecture}
L'architecture repose sur trois piliers :
\begin{itemize}
    \item \textbf{RAG (Retrieval Augmented Generation)} : Pour l'analyse de documents locaux.
    \item \textbf{APIs Médicales} : PubMed, WHO, OpenFDA.
    \item \textbf{Moteur de Décision} : Basé sur LangChain et GPT-4o-mini.
\end{itemize}

\section{Fonctionnalités}
L'assistant permet la recherche bibliographique, l'analyse de rapports PDF, et la consultation de données sur les médicaments.

\section{Conclusion}
Ce projet démontre l'efficacité des agents IA dans le secteur de la santé.

\end{document}
